
\documentclass{article}
\usepackage{colortbl}
\usepackage{makecell}
\usepackage{multirow}
\usepackage{supertabular}

\begin{document}

\newcounter{utterance}

\centering \large Interaction Transcript for game `eqbench', experiment `eq\_bench', episode 69 with BSU{-}SLIM\_\_Qwen3.5{-}9B{-}Base\_\_playpen{-}prm{-}9b{-}best\_of\_n.
\vspace{24pt}

{ \footnotesize  \setcounter{utterance}{1}
\setlength{\tabcolsep}{0pt}
\begin{supertabular}{c@{$\;$}|p{.15\linewidth}@{}p{.15\linewidth}p{.15\linewidth}p{.15\linewidth}p{.15\linewidth}p{.15\linewidth}}
    \# & \multicolumn{2}{c}{Player} && \multicolumn{2}{c}{Game Master} \\
    \hline

    \theutterance \stepcounter{utterance}  
    & & & \multicolumn{4}{p{0.6\linewidth}}{
        \cellcolor[rgb]{0.9,0.9,0.9}{
            \makecell[{{p{\linewidth}}}]{
                \texttt{\tiny{[P1$\langle$GM]}}
                \\ 
\texttt{Your task is to predict the likely emotional responses of a character in this dialogue:} \\
\\ 
\texttt{Jacques: Fancy seeing you here, Chase.} \\
\texttt{Olivia: The feeling is mutual, Duval. } \\
\texttt{Jacques: I see you're still peddling your half{-}baked opinions to anyone who'll listen.} \\
\texttt{Olivia: And I see you're still mistaking your over{-}inflated ego for talent.} \\
\texttt{Jacques: Talent? I've got more talent in my little finger than you have in your entire body.} \\
\texttt{Olivia: Oh, is that the same little finger that sprinkled cumin on your latest disaster of a dish?} \\
\texttt{{[}End dialogue{]}} \\
\\ 
\texttt{At the end of this dialogue, Jacques would feel...} \\
\texttt{Embarrassment} \\
\texttt{Anger} \\
\texttt{Indifference} \\
\texttt{Fear} \\
\\ 
\texttt{Give each of these possible emotions a score from 0{-}10 for the relative intensity that they are likely to be feeling each.} \\
\\ 
\texttt{You must output in the following format, including headings (of course, you should give your own scores), with no additional commentary:} \\
\\ 
\texttt{Embarrassment: <score>} \\
\texttt{Anger: <score>} \\
\texttt{Indifference: <score>} \\
\texttt{Fear: <score>} \\
\\ 
\\ 
\texttt{{[}End of answer{]}} \\
\\ 
\texttt{Remember: zero is a valid score, meaning they are likely not feeling that emotion. You must score at least one emotion > 0.} \\
\\ 
\texttt{Your answer:} \\
            }
        }
    }
     \\ \\

    \theutterance \stepcounter{utterance}  
    & \multicolumn{4}{p{0.6\linewidth}}{
        \cellcolor[rgb]{0.561,0.737,0.561}{
            \makecell[{{p{\linewidth}}}]{
                \texttt{\tiny{[P1$\rangle$GM]}}
                \texttt{Anger: 9} \\
\texttt{Fear: 1} \\
\texttt{Indifference: 0} \\
\texttt{Embarrassment: 2} \\
\\ 
\texttt{Jacques initiates the conflict with a confrontational tone ("Fancy seeing you here") and escalates it by accusing Olivia of spreading "half{-}baked opinions." When Olivia retorts by calling his ego a delusion and linking it specifically to a personal failure (the dish he prepared), Jacques responds with an aggressive deflection rather than addressing the specific insult. His response—"Talent? I've got more talent in my little finger..."—is a classic example of lashing out to regain the upper hand after being challenged on a weakness. The intensity of his aggression is high, driven by the need to protect his ego from being exposed. While he might feel slight embarrassment because his argument was met with a factual detail about his cooking, his primary and most dominant reaction is defensive anger. There is no evidence of fear or indifference in his aggressive posture.} \\
            }
        }
    }
    & & \\ \\

    \theutterance \stepcounter{utterance}  
    & & & \multicolumn{2}{p{0.3\linewidth}}{
        \cellcolor[rgb]{0.9,0.9,0.9}{
            \makecell[{{p{\linewidth}}}]{
                \texttt{\tiny{[GM$|$GM]}}
                \texttt{Anger: 9} \\
\texttt{Fear: 1} \\
\texttt{Indifference: 0} \\
\texttt{Embarrassment: 2} \\
            }
        }
    }
    & & \\ \\

    \theutterance \stepcounter{utterance}  
    & & & \multicolumn{2}{p{0.3\linewidth}}{
        \cellcolor[rgb]{0.9,0.9,0.9}{
            \makecell[{{p{\linewidth}}}]{
                \texttt{\tiny{[GM$|$GM]}}
                \texttt{Embarrassment: 7} \\
\texttt{Anger: 8} \\
\texttt{Indifference: 0} \\
\texttt{Fear: 0} \\
            }
        }
    }
    & & \\ \\

    \theutterance \stepcounter{utterance}  
    & & & \multicolumn{2}{p{0.3\linewidth}}{
        \cellcolor[rgb]{0.9,0.9,0.9}{
            \makecell[{{p{\linewidth}}}]{
                \texttt{\tiny{[GM$|$GM]}}
                \texttt{game\_result = LOSE} \\
            }
        }
    }
    & & \\ \\

\end{supertabular}
}

\end{document}
